\subsection{The polarized matrix algebra at a comparison arc}
\label{sec:patching-matrices}
\hypertarget{source-ordinary-patching-matrices}{}

The marking keeps the tests fixed along the chain. To recover matrix entries from
determinant data, begin with the common patch ring. It carries the universal group determinant on the
free marked cover, the ordered local characters, the ordered
Frobenius roots, the tame variables and one coefficient variable.
In its quotient \(B_0\) impose the kernel of the marking, polarization,
ordinary and inertia identities and the auxiliary trace identities against
all group elements. Its augmentation maps onto the ordered splitting
algebra of the original determinant ring. Indeed the images contain the
original determinant coefficients, ordered characters, coefficient
variable and the adjoined ordered roots, which generate that algebra.
The group and character maps give the augmentation, and each auxiliary
identity maps to the Cayley--Hamilton identity for an actual split
characteristic polynomial or to zero on tame inertia.

\begin{lemma}[Common coefficients on a marked chain]\label{lem:common-marked-coefficients}
A finite marked chain of comparison arcs admits one enlarged
finite residue field with all ordered roots and link
inclusions preserved. The arc maps are surjections to
\(k[[t]]\). At an arc with comparison ring \(B\), the marked
coefficient map extends to the completion of
\(O[[S]]_{(\varpi)}\) and identifies its residue field by
\[
 \operatorname{Frac}(k[[S]])\xrightarrow[\ S\mapsto t\ ]{\ \sim\ } k((t)).
\]
With this identification the residual coefficient map is the identity.
\end{lemma}

\begin{proof}
\proofstep{Coefficient extensions at a marked arc.}
Here are the coefficient choices implicit in completing at a point.
Enlarge the residue fields of the finite list of arcs to one finite field
\(k\), extend the coefficient DVR by the corresponding unramified extension
and adjoin \(s\), evaluated at a uniformizer \(t\) of each arc.
Write \(Q_1\supseteq P\subseteq Q_2\) for a link between adjacent curve kernels. In the finite flat coefficient extension start with a lift \(Q'_1\) of \(Q_1\). Going down and then going up give primes with contractions 
\[
 Q'_1\supseteq P'\subseteq Q'_2,\qquad P'\cap R=P,\quad Q'_2\cap R=Q_2,
\]
 where \(R\) denotes the ring before coefficient extension. Lying over applies also to a chain with just
one point.

The monic polynomial lifting the separable residue extension splits over each arc
after the chosen enlargement, so the specified prime over the curve is
the kernel of one of the evaluations. Extending each link by power series
in \(s\) retains a domain quotient and its inclusion in both adjacent
evaluation kernels. Apply the same zigzag lifting to the finite free
ordered splitting algebra. For a monic \(P(T)\) of degree \(n\), division with remainder over \(S\) gives the free module decomposition for the first root algebra
\[
 S_1=S[T_1]/(P(T_1))=\bigoplus_{j=0}^{n-1} ST_1^j .
\] In \(S_1[T]\) one has
\[
 P(T)=(T-T_1)P_1(T),\qquad \deg P_1=n-1 .
\]
Repeating with \(P_1\) gives the ordered splitting algebra, of rank \(n(n-1)\cdots1\) before selecting the prescribed factor. Select the local direct factor
with the prescribed residual ordering by lifting its residue
idempotent.

It remains finite free over the complete local ring.
Here each actual characteristic polynomial
splits with distinct roots in the fraction field of the arc, and those
roots lie in its DVR by integrality. Tensoring the fiber at the generic
point of the arc with that field gives a product indexed by the compatible orderings.
Thus every chosen prime over the arc picks an evaluation with those
roots. The maps to the arc contain the constants and send \(s\) to \(t\). Completeness supplies every series by
\[
 \sum_{i\ge0}a_i s^i\longmapsto \sum_{i\ge0}a_i t^i\qquad(a_i\in k),
\]
so the maps are surjective. Inclusion of the original Hecke ideal can be tested after
contraction. At a link with nonsemisimple residue algebra,
the field extension and the change of
group require the argument in Section~\ref{subsec:closed-link}.

Fix one such arc and let \(a:B_0\twoheadrightarrow k[[t]]\)
be the comparison map. Write \(\mathfrak p=\ker a\),
\(B=\widehat{(B_0)_{\mathfrak p}}\) and \(F=k((t))\).
The residue map \(B\to F\) extends \(a\).
Let \(O\) be the coefficient DVR with uniformizer \(\varpi\),
surjecting onto \(k\), and let \(s\in B_0\) evaluate to \(t\).
\proofstep{The coefficient DVR at the comparison point.}
The continuous substitution \(O[[S]]\to B_0\) sends \(S\) to
\(s\). Composing with \(a\) reduces the coefficients and
sends \(S\) to \(t\), so the kernel of the composition is
\((\varpi)\). For \(D_0=O[[S]]_{(\varpi)}\), this substitution gives
\[
 D_0/\varpi D_0=\operatorname{Frac}(k[[S]])
        \xrightarrow[\ S\mapsto t\ ]{\ \sim\ } k((t))=F .
\]
Every denominator outside \((\varpi)\) thus maps to a unit of \(B\); the map factors through \(D_0\). Since \(\varpi^h D_0\) maps into \(\mathfrak m_B^h\) for all \(h\) and \(B\) is complete, it extends to \(\mathcal O=\widehat{D_0}\). The identification
of its residue field with \(F\) is induced by the same
substitution. Lifts of elements of
\(k[[t]]\) can be taken in \(O[[S]]\), so their images in \(B\)
also come from \(B_0\).
\label{cor:ordinary-patching-matrices-1}

\end{proof}

\paragraph{The finite Cayley--Hamilton quotient.}
To pass from traces to an algebra acting on a module, start on the CM subgroup \(G_{\rm CM}\) of Section~\ref{sec:patching-marking}
take the closed Cayley--Hamilton quotient \(E_0\) of
\(B_0[G_{\rm CM}]\). Take the closed ideal generated by the ordered ordinary, prescribed inertia and auxiliary relations together with their images under polarization. The determinant descends through that ideal by
the trace-kernel calculation below. The resulting algebra \(E_0\) carries a determinant and an anti-involution. On a group element the latter sends \[ g\ \longmapsto\ \epsilon(g)^{1-n}(cgc^{-1})^{-1}. \] Here \(c\) is the real involution from the marking. The polarization is expressed in the frame in which the anti-involution specializes to transpose at the arc point. We use degree-\(n\)
determinants as multiplicative polynomial laws, with their
Cayley--Hamilton identity under arbitrary scalar extension.
\cite{chenevier}\label{cor:ordinary-patching-matrices-2}

\begin{auxlemma}[Detecting the full matrix span by traces]\label{lem:matrix-span-by-traces}
Let \(k\subseteq k'\) be fields, \(n\ge1\), and let \(S\subseteq M_n(k')\) satisfy \(\operatorname{tr}(XY)\in k\) for all \(X,Y\in S\).
The \(k'\)-span of \(S\) equals \(M_n(k')\) if and only if some \(X_1,\ldots,X_{n^2}\in S\) satisfy
\[
 \det\bigl(\operatorname{tr}(X_aX_b)\bigr)_{1\le a,b\le n^2}\ne0\quad\text{in }k.
\]
\end{auxlemma}

\begin{proof}
The matrix units have trace-dual basis \((E_{ji})\), since
\[
 \operatorname{tr}(E_{ij}E_{uv})=\delta_{ju}\delta_{iv}.
\]
Thus a basis selected from a spanning set has nonsingular trace matrix. Conversely a nonsingular trace matrix forces the \(n^2\) selected elements to be independent over \(k'\), so they span. The determinant belongs to \(k\), and the inclusion into \(k'\) preserves nonvanishing.
\end{proof}

\begin{proposition}[The finite polarized Cayley--Hamilton algebra]
\label{prop:marked-finite-ch-algebra}
Over the completed arc comparison \(B_0\) built from
Proposition~\ref{prop:marked-profinite-system}, the closed
Cayley--Hamilton quotient \(E_0\) of \(B_0[G_{\rm CM}]\) is finite
with its adic topology, polarized anti-involution and determinant law. When the residue characters are distinct, lifts of their diagonal
idempotents can be chosen self-adjoint. In the two-block case let
\(M\) be the mixed corner module over the coefficient ring \(B\).
At every specialization with coefficient kernel \(Q\) admitting a nonsemisimple residue quotient one has
\[
          \operatorname{Ann}_B(M)\subseteq Q .
\]
\end{proposition}

\begin{proof}
\proofstep{Finiteness and the adic topology of the quotient.}
We record the finiteness and continuity needed for the matrix
realization. The original residual determinant on \(G_{\rm CM}\)
is the direct sum of \(n\) distinct characters over the finite
residue field of \(O\). In a finite group quotient through
which they factor, let \(C\) be the corresponding residual
Cayley--Hamilton algebra. It surjects onto the product of the
\(n\) scalar factors, since distinct characters are linearly
independent. Its kernel \(J\) consists of nilpotents by the
Cayley--Hamilton equation, so \(J\) is its radical in this
finite-dimensional algebra. The diagonal units lift to
orthogonal idempotents.

For a lifted rank-one idempotent the determinant law and its corner satisfy
\[
 D(1+Xe_i)=1+X,\qquad e_iCe_i=k_0 e_i,
\]
where \(k_0\) is the residual field. These are the idempotent rank
and corner rules for determinants. Expand a product in \(J^n\) by inserting \(\sum_i e_i=1\) between
its factors. A term has indices \(i_0,\ldots,i_n\). Some
\(i_a=i_b=i\) with \(a<b\), and the corresponding subproduct lies in
\[
 e_iJ^{\,b-a}e_i\subseteq e_iJe_i=0.
\]
Every summand vanishes, hence \(J^n=0\).
In \(J/J^2\) a linear functional on the \((i,j)\)-corner
produces a cocycle for \(\chi_i\chi_j^{-1}\) by reading the
\((i,j)\) entry of a group element and dividing by \(\chi_j\).

For the chosen functional write \(x(g)\) for its value on that
entry and \(z(g)=x(g)/\chi_j(g)\). Since \(J^2=0\) in the first-order
quotient, multiplication gives
\[
 x(gh)=\chi_i(g)x(h)+x(g)\chi_j(h),\qquad
 z(gh)=z(g)+(\chi_i\chi_j^{-1})(g)z(h).
\]
This map is injective on the dual corner, since the images
of the group elements span \(C\), and \(i\ne j\) on each
nonzero radical corner. The dimension of each space of
these continuous cocycles on \(G_{\rm CM}\) is bounded by \(h\),
uniformly in the finite quotient. For example their values
at a finite topological generating set determine them,
and \(G_{\rm CM}\) is open in a finitely generated profinite group.
Products of radical lifts span every \(J^j/J^{j+1}\).
It follows that
\[
                 \dim_{k_0} C\le n+\sum_{j=1}^{n-1}(n^2 h)^j .
\]

For each coefficient truncation, adding \(u-1\) for all \(u\)
in an open normal subgroup gives an algebra quotient of a
finite group ring. Its reduction modulo the coefficient
maximal ideal has the preceding bound once that subgroup
lies in the residual character kernels. Lifting spanning
group elements and using Nakayama gives the same number
of generators over each coefficient truncation. At a fixed depth \(b\) write \(C_{U,b}\) for the resulting quotient at an open normal subgroup \(U\). If at most \(d\) generators suffice over the finite scalar ring \(R_b\), then 
\[
 U'\subseteq U\ \Longrightarrow\ C_{U',b}\twoheadrightarrow C_{U,b},\qquad |C_{U,b}|\le |R_b|^d.
\]
 After a quotient attains the maximal size, every smaller subgroup gives an isomorphism.

To choose one
set of group elements, take a spanning set after stabilization
at residue depth. Every deeper group and coefficient quotient
with that residue quotient is generated by the same elements,
again by Nakayama. For each element of the closed quotient
its expressions with these generators at finite depth define
nonempty closed sets in a fixed finite power of \(B_0\),
which is compact. Let \(C_b\) be the closed set of coefficient vectors representing the chosen element at depth \(b\). Compatibility and compactness give
\[
 C_{b+1}\subseteq C_b,\quad C_b\ne\varnothing\ (b\ge1)\quad\Longrightarrow\quad \bigcap_{b\ge1} C_b\ne\varnothing .
\] This proves finite generation in the closed quotient.
At each coefficient depth the stabilized group quotient
also shows that \(u-1\) lies in that coefficient power times
\(E_0\) on an open normal subgroup.  Take the closed quotient using the finite group and coefficient quotients. \label{cor:ordinary-patching-matrices-3}

\proofstep{Checking the trace kernel at finite precision.}
For the trace-kernel calculation first take the finite
closed Cayley--Hamilton quotient without the extra
relations. Over \(B_0\), the split residual characters
give \(n\) rank-one idempotents with scalar corners,
by the idempotent lifting argument below applied to
the original residue field. The trace radical in this quotient \(C\) is the common kernel
\[
 \bigcap_{x\in C}\ker\bigl(z\longmapsto\operatorname{tr}(xz)\bigr).
\]
Cyclicity of trace makes it a two-sided ideal. Its diagonal corners are
zero by taking \(x=e_i\). These properties hold for the
image of the ideal after any scalar extension. The
determinant of \(1+z\) for a member of this ideal is
one: use elimination in the rank-one corners. Each pivot
is one because all corrections are in the trace radical.
Off-diagonal transvections have determinant one by scaling
their entry by a conjugation over a Laurent polynomial
ring and specializing the polynomial identity at zero.

For general \(x\), let \(P_x(X)=D(X+x)\), a monic
polynomial. Cayley--Hamilton makes \(X+x\) invertible after
localizing at \(P_x\), and the trace radical is an ideal.
Thus multiplicativity in that localization gives
\[
 D(X+x+z)=D(X+x)\,
                  D(1+(X+x)^{-1}z)=D(X+x).
\]
Multiplication by a monic polynomial is injective over
any scalar ring. The equality therefore holds before
localization, and \(X=0\) proves invariance on each
trace-radical coset as a polynomial law. Thus a relation whose trace pairings
with all group elements vanish is killed as a polynomial
law, in arbitrary residue characteristic.
 The ordered ordinary and fixed inertia
generators are the relations of the ordinary pair, whose
polarized images also lie in this ideal. The closure
preserves the relations by continuity. \label{cor:ordinary-patching-matrices-4}

Extend scalars to the completed point ring. Finiteness over \(B\) identifies the resulting algebra with its completion: 
\[
 E=B\otimes_{B_0}E_0\simeq\varprojlim_b E/\mathfrak m_B^bE.
\]
 The determinant and anti-involution extend by scalar extension. There is a homomorphism
\(r:E\to M_n(F)\) given by the specialized point
representation. The point kills the defining relations,
so this homomorphism extends the given action on \(G_{\rm CM}\).
Let \(n_a\) denote the constituent dimensions. The positive cyclotomic kernel already spans the full diagonal blocks at the arc. Since \(B\to F\) is onto, 
\[
 r(E)=\bigoplus_a M_{n_a}(F)\subseteq M_n(F).
\]

The determinant of \(r(x)\) and its characteristic
polynomial are the reductions of the corresponding
polynomial law for each \(x\in E\). For \(x\in\ker r\) its residual characteristic polynomial has only the term of highest degree, so Cayley--Hamilton gives
\[
 x^n\equiv0\pmod{\mathfrak m_B E}.
\]

\paragraph{Self-adjoint rank-one idempotents.}
We need idempotents compatible with the polarization. Lift the diagonal matrix units
successively, each time in the
corner complementary to the idempotents already lifted.
For a lift \(x\) of a diagonal unit, averaging \(x\) with
\(x^*\) makes it self-adjoint, since \(2\) is a unit.
The defect \(x^2-x\) lies in \(\ker r\) and is topologically
nilpotent by the previous calculation. Write \(d=x^2-x\). Since \((2x-1)^2=1+4d\), the inverse of
\(2x-1\) exists by completeness and commutes with \(x\).
One Newton step gives
\[
      x'=x-(2x-1)^{-1}d,\qquad (x')^2-x'=(2x-1)^{-2}d^2 .
\]
Iteration squares the defect up to a unit at every step.
Its limit is a self-adjoint idempotent of the same reduction.

In successive complementary corners the resulting lifts
are orthogonal, and the final complement itself lifts the
last diagonal unit. Denote the family by \(e_i\).
For an idempotent \(e_i\) the determinant on \(1+Xe_i\) is
\((1+X)^{r_i}\) over the local ring. The leading coefficient of \((1+X)^{r_i}\) is one, so reduction
preserves its degree. Since the reduction is \(1+X\), we have \(r_i=1\).
For \(y\in e_iEe_i\), its degree-one characteristic polynomial
in this corner is \(T-\operatorname{tr}(y)\). The corner unit is
\(e_i\), so Cayley--Hamilton reads
\[
 y-\operatorname{tr}(y)e_i=0,\qquad e_iEe_i=B e_i .
\] Its trace sends \(e_i\) to \(1\).
 \cite{chenevier}\label{cor:ordinary-patching-matrices-5}

\proofstep{Lifting and matching the matrix units.}
Choose a reference index \(o(i)\) for each block, constant
throughout that block. Lift the matrix unit \(E_{o(i),i}\)
to \(z_i\in e_{o(i)}Ee_i\). The two products belong to the rank-one corners at \(i\) and
\(o(i)\). Their coefficients coincide because
\[
 \operatorname{tr}(z_i^*z_i)=\operatorname{tr}(z_i z_i^*),
 \qquad \operatorname{tr}(e_i)=\operatorname{tr}(e_{o(i)})=1.
\]
Write \(b_i\) for this common coefficient. Its residue
is \(1\). Taking its unique square root congruent to \(1\)
and dividing \(z_i\) by that root gives \(u_i\) with
\[
            u_i^*u_i=e_i,\qquad u_i u_i^*=e_{o(i)} .
\]
For the reference index use \(u_i=e_i\).  In particular
\[
      x\longmapsto (u_i x u_j^*)_{ij},
      \qquad (x_{ij})\longmapsto\sum_{i,j}u_i^*x_{ij}u_j
\]
are inverse coordinate changes. Each \(x_{ij}\) lies in the
corner with indices \(o(i),o(j)\). The matrix multiplication
rule follows by summing over the middle index and using
\(\sum_i e_i=1\). With one block this is the full algebra
\(M_n(B)\).

With two blocks with references \(p,q\), let \(M=e_pEe_q\).
Use the involution to identify \(e_qEe_p\) with \(M\).
For \(u,v\in M\), define \(\beta(u,v)\in B\) by
\(uv^*=\beta(u,v)e_p\). The involution fixes \(Be_p\),
so \(\beta\) is symmetric. Cyclicity of the trace shows that
\(u^*v=\beta(u,v)e_q\). Associativity in \(E\) gives
\(\beta(u,v)w=u\beta(v,w)\) for \(w\in M\).
Thus the rule
\[
       (b,u)(b',u')=(bb'+\beta(u,u'),\,bu'+b'u)
\]
makes \(B\oplus M\) a finite commutative \(B\)-algebra,
with an injective scalar map.

All the \(\beta(u,v)\) lie
in \(\mathfrak m_B\) because the mixed corners reduce to zero.
The coordinate changes identify \(E\) with the full space of
matrices having scalar diagonal blocks and entries in \(M\)
in the off-block rectangles, with involution transpose.
 In particular the group matrices
span the entire pattern over \(B\). Their integral
\(B_0\)-span is finite and continuous in the
\(\mathfrak m_{B_0}\)-adic sense proved for \(E_0\).
\label{cor:ordinary-patching-matrices-6}

\proofstep{The determinant law in the normalized basis.}
After extending scalars to the ring \(B\oplus M\), the determinant law on the group matrices is therefore
\[
 D(g)=\det\bigl(R(g)\bigr)
\]
in the indicated matrix realization \(R\).
One way to see the equality for any scalar algebra is to add
independent variables \(X_i\) to the diagonal and eliminate
successive pivots. Each pivot is a scalar in the enlarged
base once the previous pivots have been inverted. It is
monic in its new \(X_i\), and hence this localization is
injective. For \(i\ne j\), conjugation by \(1+(T-1)e_i\) over the
Laurent polynomial ring proves
\[
 D(1+a e_{ij})=D(1+Ta e_{ij}).
\]
Both sides are polynomials in \(T\); the polynomial ring
injects into its Laurent localization over any base.

At
\(T=0\) this gives \(D(1+a e_{ij})=1\).
On the diagonal, multiplying the idempotent factors gives
\[
  D(\operatorname{diag}(a_1,\dots,a_n))
    =\prod_i D(1+(a_i-1)e_i)=\prod_i a_i .
\]
Here the idempotent lines had rank one by the specialized
characteristic polynomial identity. Elimination now identifies
the two determinant laws before setting the \(X_i\) to zero.
The identity includes singular matrices and all coefficient
characteristics at issue.  The relation quotient imposes the ordered flag, fixed inertia and auxiliary identities. The transpose rule gives the polarization identity.
\label{cor:ordinary-patching-matrices-7}

\paragraph{Full matrix algebras and nonsemisimple residue algebras.}
The off-diagonal products determine which block pattern survives specialization. Write
\(I_M=(\beta(M,M))\). For a prime \(Q\) of \(B\),
extend to a field residue of \(B\oplus M\) over \(Q\).
Such a field exists by finiteness and the injective scalar
map. If \(M\) has nonzero image in this field, the span
contains every matrix unit by rescaling a nonzero entry
from the full mixed span. If \(M\) maps to zero, the span
is block diagonal. The mixed image is detected by its products: 
\[
 M\longmapsto0\text{ in }k'\quad\Longleftrightarrow\quad I_M\subseteq Q,
\]
 where \(k'\) is the chosen residue field. Indeed each mixed square is given by \(\beta\), and a square-zero field element vanishes.

Let \(k'/\kappa(Q)\) be the chosen residue-field extension. The group matrices have trace products in \(\kappa(Q)\) because the trace comes from \(B\). Applying Lemma~\ref{lem:matrix-span-by-traces} to the two-block calculation gives
\[
 I_M\not\subseteq Q
 \quad\Longleftrightarrow\quad
 \exists\,g_1,\ldots,g_{n^2}:\ 
 \det\!\left(\operatorname{tr}(R(g_a)R(g_b))\right)_{a,b=1}^{n^2}\ne0
 \quad\text{in }\kappa(Q).
\]
The matrices themselves are evaluated over \(k'\).
 In the one-block case the span is full.

For the other criterion suppose the residue algebra obtained
from the indicated Cayley--Hamilton kernel presents a
nonsemisimple algebra over a field with coefficient kernel
\(Q\). If \(b\in\operatorname{Ann}_B M\) lies outside \(Q\),
then all the mixed corners vanish after scalar extension
to this field. The Cayley--Hamilton quotient then becomes
the product of its two full diagonal matrix algebras.
The scalar extension statement follows by taking the span
of the kernel generators in the field group algebra, since
the group algebra surjects onto \(E\) with that defining
ideal. Any quotient of this product is semisimple, a
contradiction. This proves \(\operatorname{Ann}_B M\subset Q\). Both criteria supply the inputs for the
completed local comparison in the next subsection.
\label{cor:ordinary-patching-matrices-8}
\end{proof}
